% sections/data.tex -- Section 4 (paper_v2, cut in C4 from the full version's). The details that
% left it (cluster sizes, the McAuley data, the twin check, the title vocabulary, linking every
% ancestor) are in the appendix section on the data (appendix.tex), which becomes the supplementary
% document in B5. B2 (25 November) brought in the reviewed Vietnamese version (paper_v2_vn, L10):
% the McAuley data page as the source of the quotation, a footnote here, no appendix pointers, and
% the features as the author described them. Every number is \val{} from numbers.tex.
\section{Data}
\label{sec:data}

Amazon-Meta\footnote{\url{https://snap.stanford.edu/data/amazon-meta.html}}
\citep{leskovec2007viral} is a crawl of Amazon's product pages from the summer of 2006, with
\val{raw-records} product records for books, music CDs, DVDs and VHS videos, \val{raw-reviews}
review lines and \val{raw-categories} categories. A record gives a product's title and group, the
products that customers also bought, its category paths from the root to a leaf and its reviews.
Each review has a customer ID, a date, a rating and votes.

\subsection{Cleaning the raw data}

\paragraph{Duplicate editions.} Amazon showed the reviews of a work on the page of each of its
editions and formats, and the crawl recorded every page. The same review line therefore appears
under several products: the DVD and the VHS of one film share \val{example-shared} of their
\val{example-lines} review lines, and \val{raw-shared-pct}\% of all raw review lines also appear
under another product. A filter that only removes repeated (customer, product) pairs keeps these
copies. In the data built that way (Table~\ref{tab:data}, \emph{before}),
\val{before-clustered-pct}\% of the products sit in \val{before-clusters} clusters of editions, and
\val{before-valleak-pct}\% of the validation edges and \val{before-testleak-pct}\% of the test edges
have a copy in training: the same customer reviewing another edition of the same work. The copy
reveals the held-out edge. Such leakage leads to overestimated performance
\citep[\S3.3]{kaufman2012leakage}: a model can score well by finding copies instead of learning
what users like. The path that finds copies best is the shortest co-review path (\code{PUP}), so the
leak flatters exactly the kind of path a meta-path study compares.

We merge two products when they share at least \val{rule-edition-shared} identical review lines
(customer, date, rating, votes and helpful votes) that make up at least
\val{rule-edition-overlap-pct}\% of the smaller product's lines, ignoring lines that appear under
more than \val{rule-edition-maxrun} products. This joins \val{edition-pairs} pairs into
\val{edition-clusters} clusters, and each cluster keeps the member with the most review lines.
The later Amazon review data of \citet{mcauley2015image} have the same kind of duplicate,
``mainly due to near-identical products whose reviews Amazon merges, e.g.\ VHS and DVD versions of
the same movie'',\footnote{Quoted from the data page
\url{https://cseweb.ucsd.edu/~jmcauley/datasets/amazon/links.html} (read on 5 November 2026).} and
come with de-duplicated files. To our knowledge, no such step has been reported for Amazon-Meta.

\paragraph{Other steps.} We drop every customer account with at least
\val{rule-placeholder-repeats} reviews on a single product (\val{placeholder-ids} accounts,
\val{placeholder-lines} review lines). A review is a positive interaction when its rating is at
least \val{rule-pos-rating} (as for \val{raw-pos-pct}\% of the raw lines), unless at least
\val{rule-helpful-votes} readers voted on it and fewer than \val{rule-helpful-pct}\% found it
helpful, which drops \val{raw-unhelpful-pct}\% of the positive reviews. A user counts once per
work. A \val{rule-kcore}-core then keeps the users and products with at least \val{rule-kcore}
interactions each (\val{kcore-iterations} rounds). It is applied after merging and the account
rule so that it counts only distinct interactions.

% t1_data.tex -- T1 of the plan: the data before and after the fixes. Every number is \val{} from
% numbers.tex (make_numbers.py; sources: audit_before_fix.json and data/AMAZON/audit.json).
\begin{table}[htbp]
  \centering
  \small
  \caption{The data before and after merging duplicate editions and dropping the accounts with at
    least \val{rule-placeholder-repeats} reviews on one product. Both versions keep the positive
    reviews and the \val{rule-kcore}-core. Audits of the full data, split seed~\val{split-seed}.}
  \label{tab:data}
  \begin{tabularx}{\linewidth}{@{}Lrr@{}}
    \toprule
    & Before & After \\
    \midrule
    Products / categories & \val{before-products} / \val{before-categories}
      & \val{after-products} / \val{after-categories} \\
    Users & \val{before-users} & \val{after-users} \\
    User--product edges & \val{before-edges} & \val{after-edges} \\
    Largest user degree & \val{before-maxdeg} & \val{after-maxdeg} \\
    Duplicate-edition clusters (products in them) & \val{before-clusters} (\val{before-clustered})
      & \val{after-clusters} (\val{after-clustered}) \\
    Edges that repeat a user's interaction with another edition & \val{before-copyedges-pct}\%
      & \val{after-copyedges-pct}\% \\
    Validation / test edges with a copy in training & \val{before-valleak-pct}\% /
      \val{before-testleak-pct}\% & \val{after-valleak-pct}\% / \val{after-testleak-pct}\% \\
    Validation / test edges whose user has a training edge to a product of the same title (an
      upper bound on the duplicates left) & \val{before-valtitle-pct}\% / \val{before-testtitle-pct}\%
      & \val{after-valtitle-pct}\% / \val{after-testtitle-pct}\% \\
    \bottomrule
  \end{tabularx}
\end{table}


\paragraph{The effect.} Merging and the account rule shrink the graph from \val{before-edges} to
\val{after-edges} user--product edges and cut the test edges with a copy in training from
\val{before-testleak-pct}\% to \val{after-testleak-pct}\% (Table~\ref{tab:data}). A second check, for
products with exactly the same set of reviewers, finds \val{twin-test} test edges with such a twin
in training. Leaving them out changes no run's test Hit@10 by more than \val{twin-maxdiff}.

\subsection{The graph and the split}

A product's features are its product group and a \val{title-dim}-dimensional LSA vector of its
title. A user's features are only a \val{user-dim}-dimensional binary vector hashed from the
customer ID. They are never built from the user's reviews, because those reviews are the edges to
predict.

\paragraph{Leaf categories and near-empty paths.} Products link to their leaf categories only.
Linking every ancestor as well would put \val{cat-all-largest-pct}\% of the products into one
category. However, a single step along the hierarchy rarely leads to
another product: \code{PCCP} links a product to another one for \val{pccp-reach-dvd-pct}\% of the
DVDs and \val{pccp-reach-video-pct}\% of the videos, but for only \val{pccp-reach-book-pct}\% of the
books and \val{pccp-reach-music-pct}\% of the music. We call the six candidate paths with such a
step \emph{near-empty}: \code{PCCP}, \code{PCCPP}, \code{PCCPPP}, \code{PCCPCP}, \code{PCCPUP} and
\code{PPCCPP}. In HAN's graphs each leaves
\val{empty-isolated-min-pct}--\val{empty-isolated-max-pct}\% of the products without a neighbour,
while \code{PCCCP} (which reaches sibling categories through their parent) leaves
\val{pcccp-graph-isolated}. Books and music make up \val{group-bookmusic-pct}\% of the products and
\val{test-bookmusic-pct}\% of the test edges, so these paths come back in Section~\ref{sec:results}.

\paragraph{The split.} We split the \val{after-edges} user--product edges once at random into
\val{split-train} training, \val{split-val} validation and \val{split-test} test edges (80/10/10,
split seed~\val{split-seed}). Every graph and every walk uses the training edges only. Validation
edges guide the search and every choice made during training, and test edges only report results.
The split ignores review dates, a limitation we return to in Section~\ref{sec:conclusion}.
